\begin{pastexam}[label={px:ccu-108-swa-10}]{中正資工所甲組 108 學年度 第 11 題}
（問答）

本題接續第 10 題。 第 10 題：(1\%) What is the name of such representation when a directed graph
with edges ($u\to v$) is defined as follows?

\begin{center}
\begin{tabular}{|l|l|l|l|}\hline
$v \backslash u$ & 1 & 2 & 3\\\hline
1 & 0 & 0 & 0\\\hline
2 & 20 & 0 & -30\\\hline
3 & 40 & 0 & 0\\\hline
\end{tabular}
\end{center}

第 11 題：(1\%, 4\%) Following the above question, what does Johnson's algorithm do to such graph,
and what is the result? Write down your process.
\end{pastexam}

\begin{pastsoln}
答案：Johnson 演算法求\textbf{所有點對的最短路徑}，做法是先用 Bellman--Ford 算出 $h$，把權重改成
非負，再從每個頂點跑 Dijkstra。 結果（$\delta(u, v)$，列是起點 $u$）：
\begin{center}
\begin{tabular}{c|ccc}
$\delta(u, v)$ & $v = 1$ & $2$ & $3$ \\ \hline
$u = 1$ & $0$ & $10$ & $40$ \\
$2$ & $\infty$ & $0$ & $\infty$ \\
$3$ & $\infty$ & $-30$ & $0$ \\
\end{tabular}
\end{center}

\textbf{圖}。 依第 10 題的表（列是終點 $v$、行是起點 $u$，$0$ 表示沒有邊），邊是
$1 \to 2$：$20$、$1 \to 3$：$40$、$3 \to 2$：$-30$。 有負邊，不能直接用 Dijkstra（資結書第~25.2.5~節）；
Johnson 演算法（資結書第~25.7~節、虛擬碼~25.11）的步驟如下。

\textbf{第 1 步：加新頂點 $q$}，從 $q$ 連權重 $0$ 的邊到 $1, 2, 3$。

\textbf{第 2 步：從 $q$ 跑 Bellman--Ford}，得 $h(v) = \delta(q, v)$。 初始 $d[q] = 0$，其他 $\infty$。
每一輪依序鬆弛 $q \to 1$、$q \to 2$、$q \to 3$、$1 \to 2$、$1 \to 3$、$3 \to 2$：
\begin{itemize}
\item 第 $1$ 輪：$d[1] = d[2] = d[3] = 0$；$1 \to 2$ 得 $20$、$1 \to 3$ 得 $40$，都不更小；
$3 \to 2$ 得 $0 + (-30) = -30 < 0$，$d[2] = -30$。
\item 第 $2$、$3$ 輪：沒有任何更新。 再檢查一次所有邊也都不能再鬆弛，所以沒有負環。
\end{itemize}
得 $h(1) = 0$、$h(2) = -30$、$h(3) = 0$。

\textbf{第 3 步：改權重} $\hat{w}(u, v) = w(u, v) + h(u) - h(v)$：
\[
\hat{w}(1, 2) = 20 + 0 - (-30) = 50, \quad
\hat{w}(1, 3) = 40 + 0 - 0 = 40, \quad
\hat{w}(3, 2) = -30 + 0 - (-30) = 0.
\]
三個都非負。 由資結書引理~25.11，每條 $u$ 到 $v$ 的路徑都加上同一個數 $h(u) - h(v)$，最短路徑不變。

\textbf{第 4 步：從每個頂點用 $\hat{w}$ 跑 Dijkstra}，再換回原來的距離
$\delta(u, v) = \hat{\delta}(u, v) - h(u) + h(v)$。
\begin{itemize}
\item 從 $1$：取出 $1$，$\hat{d}[2] = 50$、$\hat{d}[3] = 40$；取出 $3$（$40$），$3 \to 2$ 得
$40 + 0 = 40 < 50$，$\hat{d}[2] = 40$；取出 $2$。 換回：$\delta(1, 2) = 40 - 0 + (-30) = 10$，
$\delta(1, 3) = 40 - 0 + 0 = 40$。 對照原圖：$1 \to 3 \to 2$ 是 $40 - 30 = 10$，比直接走 $1 \to 2$
的 $20$ 短。
\item 從 $2$：$2$ 沒有出邊，$\delta(2, 1) = \delta(2, 3) = \infty$。
\item 從 $3$：$\hat{d}[2] = 0$，換回 $\delta(3, 2) = 0 - 0 + (-30) = -30$；$1$ 走不到，$\infty$。
\end{itemize}
對角線都是 $0$。 結果就是上面的表。 這張圖只有 $3$ 個頂點，用 Floyd--Warshall 會比較省事；Johnson
演算法的好處在稀疏的大圖：$O(nm \lg n)$ 比 $\Theta(n^3)$ 快。
\end{pastsoln}
